\documentclass[11pt]{article}
\newcommand{\paso}{1}
\usepackage[T1]{fontenc}
\usepackage[utf8]{inputenc}
\usepackage[spanish,es-nodecimaldot,es-noquoting,es-lcroman]{babel}
\spanishplainpercent
\usepackage{lmodern}
\usepackage[a4paper,margin=21mm,headheight=15pt]{geometry}
\usepackage{amsmath,amssymb,bm}
\usepackage{booktabs,array}
\usepackage{tikz}
\usetikzlibrary{arrows.meta,positioning,calc}
\usepackage[most]{tcolorbox}
\usepackage{microtype}
\usepackage{fancyhdr}
\usepackage{hyperref}
\hypersetup{colorlinks=true,linkcolor=blue!45!black,urlcolor=blue!45!black}
\setlength{\parindent}{0pt}
\setlength{\parskip}{5pt}
\setlength{\emergencystretch}{2em}
\pagestyle{fancy}
\fancyhf{}
\fancyhead[L]{\small BEM del cilindro: paso \paso}
\fancyhead[R]{\small Guía matemática independiente}
\fancyfoot[C]{\thepage}
\newcommand{\vct}[1]{\bm{#1}}
\newcommand{\dd}{\,\mathrm{d}}
\newcommand{\ii}{\mathrm{i}}
\newcommand{\norm}[1]{\left\lVert#1\right\rVert}
\newcommand{\VP}{\operatorname{VP}}
\newcommand{\tr}{\operatorname{tr}}
\newtcolorbox{idea}[1]{breakable,colback=blue!3,colframe=blue!50!black,
  title={#1},fonttitle=\bfseries,boxrule=0.7pt}
\newtcolorbox{alerta}[1]{breakable,colback=orange!4,colframe=orange!65!black,
  title={#1},fonttitle=\bfseries,boxrule=0.7pt}
\newtcolorbox{resultado}[1]{breakable,colback=teal!3,colframe=teal!55!black,
  title={#1},fonttitle=\bfseries,boxrule=0.7pt}
\newcommand{\inicio}[2]{%
 {\small\color{blue!55!black}\bfseries SERIE: LOS SIETE PASOS MATEMÁTICOS DEL BEM}\par
 {\LARGE\bfseries Paso \paso. #1\par}\vspace{3mm}
 \begin{idea}{¿Qué busca este paso?}#2\end{idea}}
\newcommand{\mapa}[3]{%
 \begin{center}
 \begin{tikzpicture}[node distance=4mm,
   caja/.style={draw=blue!50!black,fill=blue!3,rounded corners=2pt,
     text width=13cm,align=center,inner sep=6pt},
   >={Stealth[length=2mm]}]
 \node[caja] (a){#1};
 \node[caja,below=of a] (b){#2};
 \node[caja,below=of b] (c){#3};
 \draw[->,thick,blue!60!black] (a)--(b);
 \draw[->,thick,blue!60!black] (b)--(c);
 \end{tikzpicture}
 \end{center}}
\newcommand{\base}{%
 \smallskip{\footnotesize\color{black!65}
 Modelo de referencia: cilindro sólido, material homogéneo e isótropo,
 deformación plana, pequeñas deformaciones, sin fuerzas de volumen;
 amplitudes complejas con \(e^{+\ii\omega t}\), \(\omega>0\).}\par}
\newcommand{\fuentes}{%
 \par\medskip\begingroup\footnotesize
 \textbf{Referencia y continuación.}
 Este documento desarrolla una etapa del esquema
 \texttt{esquema\_ecuaciones\_bem.tex}. Las ecuaciones y los ejemplos se
 contrastaron con \texttt{bem\_cilindro.py}, ambos en la carpeta superior.
 La guía general y sus deducciones ampliadas se encuentran en
 \texttt{../guia\_bem/guia\_bem.pdf}. No se modifica la implementación.
 \par Para recompilar este documento, mantener \texttt{estilo.tex} en
 la misma carpeta y ejecutar \texttt{pdflatex} sobre su archivo
 \texttt{.tex}, desde esa carpeta.\par\endgroup}

\begin{document}
\inicio{El problema físico}{Traducir el cilindro sometido a presión en un
problema matemático bien definido: qué región se estudia, qué campo se busca,
qué ecuación cumple y qué se conoce en el borde. Todavía no se integra ni se
construye ninguna matriz.}
\base
\section*{1. La pregunta de ingeniería}
Se conoce la presión que comprime el contorno. Se busca cuánto se desplaza el
material cuando esa presión oscila con frecuencia angular \(\omega\).
El resultado no es necesariamente una contracción instantánea en fase con la
carga: la inercia y el amortiguamiento producen amplitud y desfase.

\begin{center}
\begin{tikzpicture}[scale=0.85,>={Stealth[length=2mm]}]
 \fill[blue!5] (0,0) circle(1.6);
 \draw[thick,blue!65!black] (0,0) circle(1.6);
 \foreach \a in {0,90,180,270}
   \draw[->,thick,red!65!black] (\a:2.2)--(\a:1.65);
 \draw[->,thick,teal!65!black] (1.13,1.13)--(1.85,1.85)
   node[right]{$\vct n$ exterior};
 \draw[->] (0,0)--(1.6,0) node[midway,below]{$R$};
 \node at (-0.45,0.5){$\Omega$};
 \node at (-1.6,1.5){$\Gamma$};
 \node[red!65!black] at (0,-2.65){$\vct t=-p\vct n$};
 \node[align=left,text width=7cm] at (6,0){
   $\Omega=\{(x,y):x^2+y^2<R^2\}$\\[2mm]
   $\Gamma=\partial\Omega$\\[2mm]
   El área contiene material.\\
   No se modela una cavidad ni el exterior.};
\end{tikzpicture}
\end{center}

\mapa{Presión conocida \(p\), geometría \(R\), material y frecuencia \(\omega\)}
{Equilibrio dinámico + deformación + ley constitutiva}
{Campo buscado \(\vct u(x,y;\omega)\), en particular su valor radial en \(\Gamma\)}

\section*{2. Qué entra y qué sale}
\begin{tabular}{p{0.20\linewidth}p{0.70\linewidth}}
\toprule
Símbolo & Significado y unidad\\
\midrule
\(R,p,\rho\) & Radio [m], amplitud de presión [Pa], densidad [kg/m\(^{3}\)].\\
\(\vct u\) & Desplazamiento complejo \((u_x,u_y)\), en metros.\\
\(\boldsymbol\varepsilon\) & Deformación infinitesimal, adimensional.\\
\(\boldsymbol\sigma,\vct t\) & Tensor de esfuerzos y vector de tracción [Pa].\\
\(\omega\) & Frecuencia angular [rad/s]; no es frecuencia en Hz: \(\omega=2\pi f\).\\
\bottomrule
\end{tabular}

\newpage
\section*{3. Del desplazamiento a la ecuación de campo}
\subsection*{3.1 La cinemática: cómo se deforma el material}
Un desplazamiento constante mueve el sólido pero no lo deforma. Lo que genera
deformación son las variaciones espaciales de \(\vct u\):
\[
 \varepsilon_{ij}=\frac12(u_{i,j}+u_{j,i}),\qquad
 \varepsilon_{xx}=u_{x,x},\quad
 \varepsilon_{yy}=u_{y,y},\quad
 \varepsilon_{xy}=\frac12(u_{x,y}+u_{y,x}).
\]
La notación \(u_{i,j}\) significa \(\partial u_i/\partial x_j\).
No confundir \(\varepsilon_{xy}\) con la deformación cortante ingenieril
\(\gamma_{xy}=2\varepsilon_{xy}\).

\subsection*{3.2 La constitución: cómo la deformación genera esfuerzo}
Para elasticidad lineal isótropa, usando los módulos complejos del paso 2:
\[
 \boxed{\sigma_{ij}=\lambda^* \varepsilon_{kk}\delta_{ij}
                  +2\mu^*\varepsilon_{ij}.}
\]
Un índice repetido implica suma. Aquí \(\delta_{ij}\) vale uno si \(i=j\)
y cero en caso contrario. En deformación plana,
\(\varepsilon_{zz}=\varepsilon_{xz}=\varepsilon_{yz}=0\), pero
\[
 \sigma_{zz}=\lambda^*(\varepsilon_{xx}+\varepsilon_{yy})
\]
puede ser distinto de cero: la restricción axial produce esfuerzo axial.

\subsection*{3.3 El equilibrio: de Newton a Navier}
Sin fuerzas de volumen, la segunda ley de Newton por unidad de volumen es
\[
 \sigma_{ij,j}=\rho\,\ddot u_i.
\]
Si \(u_{\text{fís},i}=\operatorname{Re}\{u_i e^{\ii\omega t}\}\), entonces
\(\ddot u_{\text{fís},i}=\operatorname{Re}\{-\omega^2u_i e^{\ii\omega t}\}\).
Por tanto,
\[
 \boxed{\sigma_{ij,j}+\rho\omega^2u_i=0.}
\]
Al sustituir Hooke y agrupar derivadas, para material homogéneo:
\[
 \boxed{\mu^*\nabla^2\vct u+
 (\lambda^*+\mu^*)\nabla(\nabla\!\cdot\vct u)
 +\rho\omega^2\vct u=\vct0.}
\]
\begin{alerta}{El signo de la inercia no se adivina}
El término \(+\rho\omega^2\vct u\) aparece después de llevar
\(-\rho\omega^2\vct u\) al lado izquierdo. El amortiguamiento está en los
módulos; no se añade además un término viscoso que el modelo no contiene.
\end{alerta}

\newpage
\section*{4. De la presión escalar a la condición vectorial}
La tracción sobre una superficie orientada es
\[
 t_i=\sigma_{ij}n_j,\qquad \boxed{\vct t=-p\vct n\quad\text{en }\Gamma.}
\]
En el círculo, \(\vct x=(R\cos\theta,R\sin\theta)\) y
\[
 \vct n=(\cos\theta,\sin\theta),\qquad
 \vct t=(-p\cos\theta,-p\sin\theta).
\]
Para \(p=100\) Pa:
\begin{center}
\begin{tabular}{ccc}
\toprule
Punto & Normal exterior & Tracción [Pa]\\
\midrule
\((R,0)\)&\((1,0)\)&\((-100,0)\)\\
\((0,R)\)&\((0,1)\)&\((0,-100)\)\\
\((-R,0)\)&\((-1,0)\)&\((100,0)\)\\
\((0,-R)\)&\((0,-1)\)&\((0,100)\)\\
\bottomrule
\end{tabular}
\end{center}
La presión es uniforme, pero las componentes cartesianas de la tracción no
son constantes: cambian con la orientación de la normal.

\section*{5. Por qué el problema es bidimensional}
Se considera una respuesta independiente de \(z\), con desplazamiento axial
nulo. Esto representa la idealización de deformación plana de un cilindro
largo axialmente restringido, o una región suficientemente alejada de sus
extremos bajo esa hipótesis. No basta que el cilindro sea largo para
justificar cualquier restricción axial.

\begin{alerta}{Deformación plana no es esfuerzo plano}
Aquí \(\varepsilon_{zz}=0\), no \(\sigma_{zz}=0\). Si los extremos y las
condiciones axiales permiten otra respuesta, puede corresponder una
formulación de deformación plana generalizada u otro modelo. No cambiar
\(\lambda^*\) por una fórmula de esfuerzo plano sin cambiar el problema.
\end{alerta}

\section*{6. Unidades como comprobación}
\[
 [\nabla\!\cdot\boldsymbol\sigma]=\mathrm{Pa/m}
 =\mathrm{N/m^3},\qquad
 [\rho\omega^2\vct u]=\mathrm{kg/m^3}\,\mathrm{s^{-2}}\,\mathrm m
 =\mathrm{N/m^3}.
\]
Ambos términos son fuerzas por unidad de volumen. La presión, en cambio,
es fuerza por unidad de área y actúa en el borde.

\newpage
\section*{7. Qué implica la simetría y qué no}
Para geometría circular, material isótropo y presión uniforme, la solución
física es axisimétrica:
\[
 \vct u(r,\theta)=u_r(r)\vct e_r,\qquad u_\theta=0.
\]
Eso sirve para comprobar el resultado, pero el BEM empleado no reduce
previamente el sistema a una sola incógnita radial: conserva las dos
componentes cartesianas en todos los nodos.

\begin{resultado}{Ejemplo estático que permite revisar los signos}
En el límite estático, un campo \(u_r(r)=Ar\) tiene
\(\varepsilon_{rr}=\varepsilon_{\theta\theta}=A\). Por Hooke:
\[
 \sigma_{rr}=2(\lambda^*+\mu^*)A=-p,
 \qquad
 \boxed{u_r(R)=-\frac{pR}{2(\lambda^*+\mu^*)}.}
\]
Sin amortiguamiento, con \(R=6\) m y
\(\lambda=\mu=10^6\) Pa, resulta
\(u_r(R)=-1.5\times10^{-4}\) m. El signo negativo indica contracción.
Este es un control del modelo, no una evaluación directa de los núcleos
dinámicos en \(\omega=0\).
\end{resultado}

\section*{8. Aclaraciones antes de continuar}
\begin{itemize}
 \item Una amplitud compleja no es una posición física compleja. La
 posición física se obtiene tomando la parte real de la respuesta temporal.
 \item La hipótesis lineal exige desplazamientos y gradientes de
 desplazamiento suficientemente pequeños; no incluye grandes deformaciones.
 \item Se conoce toda la tracción del borde. No se prescribe simultáneamente
 todo el desplazamiento del mismo borde.
 \item La solución fundamental usada después presupone homogeneidad e
 isotropía dentro del dominio. Un material variable requiere otro tratamiento.
\end{itemize}

\mapa{Modelo físico definido: \(\Omega,\Gamma,\rho,p,\omega\)}
{Paso 2: obtener \(\lambda^*,\mu^*,c_p,c_s\)}
{Paso 3: sustituir la ecuación de campo por una identidad en \(\Gamma\)}
\fuentes
\end{document}
